\documentclass[11pt,a4paper]{article}
\usepackage[T1]{fontenc}
\usepackage[utf8]{inputenc}
\usepackage[polish]{babel}
\usepackage{amsmath,amssymb,amsthm,geometry,hyperref}
\geometry{margin=25mm}
\newtheorem{theorem}{Twierdzenie}
\newtheorem{corollary}{Wniosek}
\newcommand{\Tr}{\operatorname{Tr}}
\newcommand{\Rea}{\operatorname{Re}}
\newcommand{\Dtr}{D_{\rm tr}}
\title{FIN: równoważność dynamiki Hartreego i wymuszona korelacja strict}
\author{}
\date{}
\begin{document}
\maketitle

\begin{abstract}
Sama zgodność czystej dynamiki projektowej nie wyznacza oddziaływania
dwuciałowego: dla 12 etykiet pozostawia 4357 rzeczywistych parametrów.
Zgodność dla wszystkich stanów mieszanych pozostawia tylko dwa, a
znajomość pełnej mapy źródłowej usuwa tę swobodę. Pomimo tak dużej
nieoznaczoności mikroskopowej istnieje element komutatora strict,
którego nie zmienia żaden z niewidocznych parametrów. Wyklucza on
stacjonarne produkty stanów strict niezależnie od ich rzędu.
Dla dwóch nośników i zadanej wartości $\gamma=1/20$ każdy dopuszczalny
stan stacjonarny o tych marginałach spełnia
$\Dtr(R,C\otimes E)>0.0056$. Jest to jednolita dolna granica korelacji
w całej zadeklarowanej klasie, nie tylko niepowodzenie jednej konstrukcji.
Nie stanowi ona pomiaru, konieczności splątania ani domknięcia FIN
jako fundamentalnej teorii fizycznej.
\end{abstract}

\section{Klasa modeli i rozdzielenie przesłanek}

Punktem wyjścia jest mapa $T(\rho)=\Pi\Rea\rho$, gdzie $\Pi$ usuwa
przekątną w ustalonej bazie etykiet. Jest to źródło z badanego
sprzężenia uczącego, którego granica szybkiego śledzenia prowadzi
do $\dot\rho=(i/\gamma)[T(\rho),\rho]$.
\footnote{Źródła repozytoryjne: \texttt{The FIN Kernel as an Unknown Mat.md},
§4.1; \texttt{fin\_projected\_learning/REPORT.md}, §28;
\texttt{FIN\_Quantum\_Source\_Correlation\_and\_Programming\_Report.tex}.
Obecny wynik nie zastępuje pełnego równania skończonej szybkości uczenia
jego granicą bez zaznaczenia tej zmiany.}

Niech $S$ zamienia dwa czynniki, $P_s=(I+S)/2$, $P_a=(I-S)/2$,
$|\Omega\rangle=\sum_i|ii\rangle$ oraz
$D=\sum_i|ii\rangle\langle ii|$. Kanoniczne podniesienie mapy źródłowej to
\begin{equation}\label{eq:V}
 V=\frac{|\Omega\rangle\langle\Omega|+S-2D}{2},\qquad
 L_V(\rho)=\Tr_2[V(I\otimes\rho)]=T(\rho).
\end{equation}
Porównujemy stałe hermitowskie oddziaływania dwuciałowe $H$.
Warunek $SHS=H$ oznacza tu symetrię wymiany nośników, nie dodatkowe
twierdzenie o odwróceniu czasu lub fizycznej statystyce cząstek.
Oznaczamy
$F_H(\rho)=i[L_H(\rho),\rho]$.

Wszystkie klasyfikacje dotyczą przepływu projektorów lub macierzy
gęstości. Nie ustalają kontrolowanej fazy bramki, początku zegara,
jednostek ani kanonicznego selektora. Wymiar nośnika $n=12$ nie jest
utożsamiany z liczbą fizycznych cząstek.

\begin{theorem}[Trzy poziomy identyfikacji]\label{thm:classes}
Dla oddziaływań symetrycznych względem wymiany:
\begin{enumerate}
\item Jeśli $L_H(P)=T(P)$ dla każdego projektora rzędu jeden, to $H=V$.
\item Jeśli $F_H(\rho)=F_V(\rho)$ dla każdej macierzy gęstości, to
$H=V+aS+bI$, gdzie $a,b\in\mathbb R$.
\item Jeśli równość przepływów jest wymagana tylko dla projektorów
rzędu jeden, to dokładnie
\begin{equation}\label{eq:pureclass}
 H=V+cP_s+B_a,\qquad B_a=P_aB_aP_a=B_a^*,\quad c\in\mathbb R.
\end{equation}
\end{enumerate}
\end{theorem}

\begin{proof}
Projektory rzędu jeden rozpinają przestrzeń hermitowską, więc pierwsza
teza wynika z liniowości i niezdegenerowanego parowania śladu częściowego.
Dla drugiej rozważmy różnicę map $L$. Wielomianowa tożsamość
$[L(\rho),\rho]=0$ na wnętrzu zbioru stanów rozszerza się najpierw
na hiperpłaszczyznę śladu jeden, a przez jednorodność i ciągłość na
wszystkie macierze hermitowskie. Standardowa klasyfikacja map komutujących
daje $L(X)=aX+f(X)I$.
\footnote{M.~Brešar, Š.~Špenko, \emph{Functional identities in one variable},
wersja autorska, §3, w szczególności Lemmat 3.2,
\href{https://users.fmf.uni-lj.si/bresar/documents/FIonevar.pdf}{tekst źródłowy}.
Zasada klasyfikacji map komutujących jest znana; nie jest nowym
ogólnym twierdzeniem tego raportu.}
Można to również otrzymać przez polaryzację na diagonalnych macierzach
i rzeczywistych oraz urojonych jednostkach pozadiagonalnych: współczynnik
części nieskalarnej musi być jeden dla wszystkich par indeksów.
Symetria wymiany czyni $L$ samosprzężoną w parowaniu Hilberta--Schmidta,
co wymusza $f(X)=b\Tr X$. Reprezentantami są $S$ i $I$.

Dla trzeciej tezy niech $Z=H-V$. Na sferze jednostkowej funkcja
$\mathcal E(\psi)=\langle\psi\psi|Z|\psi\psi\rangle$ ma pochodną
$4\Rea\langle d\psi|L_Z(P)|\psi\rangle$.
Zerowy przepływ projektowy oznacza, że $L_Z(P)\psi$ jest proporcjonalne
do $\psi$, więc pochodna w każdym kierunku stycznym znika.
Funkcja $\mathcal E$ jest stała, równa $c$. Jednorodność i polaryzacja
wielomianu w zmiennych zespolonych dają $P_s ZP_s=cP_s$.
Blok antysymetryczny jest dowolny. Odwrotnie: anihiluje on
$\psi\otimes\psi$, a $cP_s$ dodaje tylko wspólną fazę.
\end{proof}

\begin{center}
\begin{tabular}{p{6.1cm}p{5.0cm}r}
Przesłanka & Pozostająca różnica $H-V$ & Wymiar dla $n=12$\\\hline
Pełna mapa źródłowa na stanach czystych & $0$ & 0\\
Przepływ wszystkich stanów mieszanych & $aS+bI$ & 2\\
Tylko czysty przepływ projektowy & $cP_s+B_a$ & 4357
\end{tabular}
\end{center}
Ostatni wymiar wynosi $1+66^2$; po odjęciu globalnego przesunięcia
energii pozostaje 4356. Są to parametry matematycznej przestrzeni
hamiltonianów, nie liczba fizycznych pól. Znajomości rzeczywistego
źródła $T$ nie wolno zastępować słabszą znajomością samego przepływu.

\section{Rzeczywiste wyjątki: dlaczego poprzedniej tezy nie wolno rozszerzać bez dowodu}

Dla $V_a=V+aS$ lokalny komutant jest skalarny przy $n\ge3$,
z jednym wyjątkiem: $a=-1/2$. Wtedy
$[V_a,h\otimes I+I\otimes k]=0$ zachodzi dokładnie dla diagonalnych
$h,k$ z $h+k=cI$. Wynika to z rozdzielenia czterech pasm:
$\Omega$, pozadiagonalnego symetrycznego, diagonalnego bezśladowego
i antysymetrycznego. Ich jedyne możliwe zbieżności są przy
$a=0,-1/2,-n/4$; test na diagonalnych i pozadiagonalnych wektorach
usuwa dodatkową swobodę w dwóch pozostałych przypadkach.

Przykładem pełnorządowej niejednorodnej równowagi iloczynowej jest
\[
 C=\operatorname{diag}(1/2,1/3,1/6),\qquad
 E=\operatorname{diag}(2,3,6)/11
\]
dla $V-S/2$. Na grafie dwudzielnym można naprzemiennie przypisać
$C$ i jego znormalizowaną odwrotność. Cykl nieparzysty wymusza
jednorodność. Pola lokalne muszą komutować z przypisanymi stanami.
Przykład obala rozszerzenie wcześniejszego kanonicznego twierdzenia
na całą nową klasę bez analizy. Nie realizuje jednak niezerowego jądra:
$T(C)=T(E)=0$.

Jeszcze słabsza przesłanka czystej dynamiki dopuszcza także niezerowy,
lecz rzadki przykład
\begin{equation}\label{eq:matchingexample}
 C=\frac15\begin{pmatrix}2&1&0\\1&2&0\\0&0&1\end{pmatrix},\qquad
 E=\frac17\begin{pmatrix}2&-1&0\\-1&2&0\\0&0&3\end{pmatrix}.
\end{equation}
W bazie antysymetrycznej $(|01\rangle-|10\rangle)/\sqrt2$,
$(|02\rangle-|20\rangle)/\sqrt2$,
$(|12\rangle-|21\rangle)/\sqrt2$ przyjmijmy blok
\[
 H_a=\begin{pmatrix}-1/2&0&0\\0&0&1/2\\0&1/2&0\end{pmatrix},
 \qquad H_s=P_sVP_s.
\]
Dokładnie $[H,C\otimes E]=0$. Jądro ma tylko jedną krawędź,
a znak krawędzi partnera jest przeciwny. Dowolne obrócenie bazy
przy zachowaniu tego samego $T$ nie jest legalnym obejściem:
normalne porządkowanie wybiera ustaloną algebrę diagonalną.

\section{Przeszkody niezależne od pełnego rzędu lub wybranego bloku}

\begin{theorem}[Identyczne czynniki, dowolny rząd]\label{thm:rankfree}
Dla $n\ge3$, dowolnego $H$ w klasie \eqref{eq:pureclass} i dowolnego
stanu $C$,
\[
 [H,C\otimes C]=0\quad\Longleftrightarrow\quad C=I/n.
\]
\end{theorem}

\begin{proof}
Żadna niezerowa właściwa podprzestrzeń $M\subset\mathbb C^n$ nie ma
$\operatorname{Sym}^2M$ niezmienniczej względem $V_s=P_sVP_s$.
Istotnie, dla $\psi\in M$ działanie $V_s$ po odjęciu
$\psi\otimes\psi/2$ odpowiada diagonalnej macierzy
$\operatorname{diag}(q/2-\psi_i^2)$, $q=\sum_i\psi_i^2$,
której obraz musiałby leżeć w $M$.
Jeżeli $e_j\in M$, wybór $\psi=e_j$ daje macierz odwracalną
i wymusza $M=\mathbb C^n$. W przeciwnym razie żadna baza współrzędna
nie należy do $M$, więc wszystkie współczynniki diagonalne muszą
znikać. Równości $\psi_i^2=q/2$ dają $(n-2)q=0$, a następnie
$\psi=0$, sprzeczność. Argument obejmuje zespolone podprzestrzenie.

Stacjonarność produktu wymuszałaby niezmienniczość jego nośnika
$M\otimes M$, zatem także $\operatorname{Sym}^2M$.
Wynika stąd pełny rząd $C$. Wtedy logarytm produktu daje pole
$h\otimes I+I\otimes h$, $h=\log C$. Jego komutacja z $V_s$
wymusza $h$ skalarne: izolowany wektor $\Omega$ ustala część
rzeczywistą, a rozdzielenie diagonalnego i pozadiagonalnego pasma
usuwa część urojoną. Stąd $C=I/n$.
\end{proof}

Wyjątek $n=2$ jest rzeczywisty: dla $V=X\otimes X/2$ istnieją
niejednorodne identyczne produkty stacjonarne.
Dla nierównych czynników pełnego rzędu w klasie \eqref{eq:pureclass}
można natomiast dowieść warunku koniecznego
\begin{equation}\label{eq:matching}
 E=C^{-1}/\Tr(C^{-1}),\qquad
 \text{pozadiagonalny graf $C$ ma stopień najwyżej jeden}.
\end{equation}
W skrócie: logarytmy dają $h+k=cI$. Dla
$L=P_s(h\otimes I-I\otimes h)P_a$ komutacja bloków i hermitowskość
wymuszają $[V_s,LL^*]=0$. Projektor $D$ jest sumą projektorów
spektralnych $V_s$, więc także z nim komutuje $LL^*$.
W obrazie macierzy diagonalnej współczynnik $d_k$ w elemencie
pozadiagonalnym $ij$, dla różnych $i,j,k$, wynosi $-2h_{ik}h_{jk}$.
Musi zniknąć. Każda kolumna $h$ ma najwyżej jedną krawędź,
a wykładnik zachowuje ten podział na bloki jedno- i dwuwymiarowe.
To dowodzi \eqref{eq:matching}, także dla zespolonego $C$.

Gęsty strict oraz kanoniczny cykl legacy nie spełniają tego warunku.
Ich co najmniej dwie krawędzie przy jednym węźle są niezerowe.
Jest to osobne zastosowanie strukturalne do obu referencji, bez
zamiany jąder i bez transferu ich ról fizycznych.

\section{Wspólne mody strict usuwają także ograniczenie rzędu partnerów}

Niech $K_\star$ będzie dostarczoną macierzą strict na $C_{12}$:
\[
 (K_\star)_{ij}=\frac{\cos((743d_{ij}+650)/4000)}{1+d_{ij}^{9/5}},
 \quad (K_\star)_{ii}=0.
\]
Oznaczmy przez $u$ jednorodny wektor Perrona i przez
$v_i=(-1)^i/\sqrt{12}$ mod naprzemienny. Obie przestrzenie własne
są proste. Ich wartości własne to $\lambda_0=s$ i $\lambda_6$.
Własności te mają oddzielne dokładne przedziały Fourierowskie.
\footnote{Źródło: \texttt{fin\_replication\_consistency/certify.py} oraz
\texttt{fin\_projected\_learning/research.py}, funkcja
\texttt{certify\_strict\_spectrum}. Przedziały są przeliczane w
certyfikacie bieżącego wyniku, a nie zastępowane zaokrąglonym widmem.}

Nie zakładamy pełnego rzędu. Zakładamy jedynie, że $C,E$ są stanami
indywidualnie stacjonarnymi i samospójnymi z tym samym strict:
\begin{equation}\label{eq:states}
 T(C)=\gamma_CK_\star,\quad[K_\star,C]=0,\qquad
 T(E)=\gamma_EK_\star,\quad[K_\star,E]=0,
 \quad\gamma_C,\gamma_E>0.
\end{equation}
Spójność grafu wymusza
$\Rea C=I/n+\gamma_CK_\star$, analogicznie dla $E$.
Urojona część jest rzeczywistą macierzą antysymetryczną komutującą
z $K_\star$; na każdej prostej rzeczywistej przestrzeni własnej musi
znikać. Zatem
\[
 Cu=c_0u,\ Cv=c_6v,\quad c_j=1/n+\gamma_C\lambda_j,
 \qquad Eu=e_0u,\ Ev=e_6v,\quad e_j=1/n+\gamma_E\lambda_j.
\]
Mamy $c_0>c_6\ge0$, $e_0>e_6\ge0$, a więc
$\Delta=c_0e_0-c_6e_6>0$.

\begin{theorem}[Niezmienny element komutatora]\label{thm:entry}
Dla $\alpha=u\otimes u$, $\beta=v\otimes v$ i
$\kappa=(n-2)/(2n)=5/12$, każdy $H$ z \eqref{eq:pureclass} spełnia
\begin{equation}\label{eq:entry}
 \langle\beta|[H,C\otimes E]|\alpha\rangle=\kappa\Delta>0.
\end{equation}
Żaden produkt stanów \eqref{eq:states} nie jest zatem równowagą
w całej tej klasie, bez względu na ich rzędy i chiralności.
\end{theorem}

\begin{proof}
Dokładnie $V\alpha=\alpha/2+\kappa\Omega$ i
$V\beta=\beta/2+\kappa\Omega$.
Blok antysymetryczny anihiluje oba wektory, a przesunięcie $cP_s$
nie zmienia badanego komutatora. Wektory $\alpha,\beta$ są własne
dla $C\otimes E$, z wartościami $c_0e_0,c_6e_6$.
To daje \eqref{eq:entry} bez założenia odwracalności stanów.
\end{proof}

Wynik obejmuje między innymi stany minimalnego rzędu sześć z przeciwnymi
chiralnościami: ich wspólny mod Perrona nie znika.
Istnieje także wersja grafowa. W produkcie wielu marginałów
\eqref{eq:states} porównujemy $u^{\otimes N}$ z wektorem zawierającym
$v$ dokładnie na końcach wybranej krawędzi. Tylko ta krawędź może
połączyć te wektory. Dowolne pola jednociałowe i inne krawędzie nie
kasują elementu \eqref{eq:entry}, pomnożonego przez sprzężenie oraz
dodatnie czynniki $c_0$ pozostałych nośników. Warunek niezerowego
sprzężenia kanonicznego na tej krawędzi jest istotny.

\section{Ilościowa dolna granica korelacji}

\begin{theorem}[Granica niezależna od niewidocznego bloku]\label{thm:floor}
Niech $H$ należy do \eqref{eq:pureclass}, a $R$ będzie dowolnym
stacjonarnym stanem dwuciałowym o marginałach $C,E$ z \eqref{eq:states}.
Wtedy
\begin{equation}\label{eq:floor}
 \Dtr(R,C\otimes E):=\tfrac12\|R-C\otimes E\|_1
 \ge\frac{\Delta}{\sqrt{2(n-2)}}.
\end{equation}
\end{theorem}

\begin{proof}
Niech $\chi=R-C\otimes E$. Jej oba ślady częściowe są zerowe.
Dla $d=\alpha-\beta$ zdefiniujmy
\begin{align}
 B_0&=\frac\kappa2\bigl(|d\rangle\langle\Omega|
                         +|\Omega\rangle\langle d|\bigr),\\
 L&=(|u\rangle\langle u|-|v\rangle\langle v|)\otimes I
       +I\otimes(|u\rangle\langle u|-|v\rangle\langle v|),\qquad
 B=B_0-\kappa L/2.
\end{align}
Stacjonarność i \eqref{eq:entry} dają
$\Tr(\chi B_0)=-\kappa\Delta$.
Odejmowanie $L$ jest dozwolone właśnie dzięki zerowym marginałom
$\chi$, więc $\Tr(\chi B)=-\kappa\Delta$.

Norma $B$ jest dokładnie $\kappa\sqrt{(n-2)/2}$.
Na rozpiętości $d$ i $\Omega-\alpha-\beta$ operator ma dwie przeciwne
wartości własne tej wielkości; na sektorach z jednym czynnikiem $u$
lub $v$ i jednym ortogonalnym ma wartości $\pm\kappa/2$.
Pozostałe wartości są zerowe. Nierówność Höldera daje
$\|\chi\|_1\ge\sqrt{2/(n-2)}\Delta$, czyli \eqref{eq:floor}.
\end{proof}

Dla $\gamma_C=\gamma_E=1/20$ dokładny przedział daje
\[
 \Delta\ge
 \frac{302968216165207140298707455443}
      {12000000000000000000000000000000}
 >\frac{63}{2500}.
\]
Ponieważ $\sqrt{20}<9/2$, dostajemy ścisły certyfikat
\begin{equation}\label{eq:number}
 \Dtr(R,C\otimes E)>\frac7{1250}=0.0056.
\end{equation}
Standardowa nierówność Pinskara, dla logarytmów naturalnych,
daje następnie
\[
 I(A:B)_R=D(R\|C\otimes E)
 \ge2\Dtr(R,C\otimes E)^2>\frac{49}{781250}\ \text{nat}.
\]
\footnote{E.~A.~Carlen, E.~H.~Lieb,
\emph{Remainder Terms for Some Quantum Entropy Inequalities} (2014),
równania (2.6)--(2.7),
\href{https://arxiv.org/abs/1402.3840}{arXiv:1402.3840}.
Nierówność Pinskara jest narzędziem literaturowym, nie nowym twierdzeniem FIN.}

Granica jest niezależna od dowolnego $B_a$, od $c$ i od wspólnej
niezerowej skali czasowej całego $H$. Nie gwarantuje, że dla każdego
$B_a$ istnieje odpowiedni stan $R$. Klasa nie jest jednak pusta:
wcześniejsza konstrukcja antysymetryczna dla kanonicznego $V$ i
$C=E=I/12+K_\star/20$ dostarcza stanu stacjonarnego o tych marginałach.

\paragraph{Granica zakresu.}
Wersja ilościowa dotyczy dwóch nośników i hamiltonianu w klasie
\eqref{eq:pureclass}, ewentualnie pomnożonego przez jedną skalę.
Dowolne dodatkowe pola lokalne spoza tej klasy wymagają nowego
oszacowania. Grafowe wykluczenie produktu dopuszcza takie pola,
ale nie przenosi na nie automatycznie liczby \eqref{eq:number}.
Nie jest to granica splątania, procent fizycznej energii ani czułość
aparatury. Dotyczy korelacji względem iloczynu zadanych marginałów.

\section{Kontrole, status wyników i następny kierunek}

\paragraph{Co udowodniono.}
Wyznaczono dokładną różnicę między identyfikacją pełnego źródła,
mieszanego przepływu i czystego przepływu projektowego. W ostatniej,
najsłabszej klasie wyprowadzono przeszkodę niezależną od rzędu stanów
strict i jednolitą dodatnią dolną granicę wymaganych korelacji.
Nie wystarcza więc dobór innego, niewidocznego bloku hamiltonianu,
aby uzyskać niezależne stacjonarne nośniki tego samego strict.

\paragraph{Co obalono lub zawężono.}
Nie wolno utożsamiać czystej dynamiki z jej mieszaną ekstrapolacją.
Nie wolno też rozszerzać kanonicznej tezy o jedynym maksymalnie
mieszanym produkcie na dowolny zmieniony lift: przykład odwrotnych
rozkładów i przykład \eqref{eq:matchingexample} są rzeczywistymi wyjątkami.
Nie realizują one jednak gęstego, samospójnego strict.
Wyjątek $n=2$, zerowe kodowanie i zmiana bazy bez zmiany $T$ zostały
oddzielnie uwzględnione.

\paragraph{Dowody a obliczenia.}
Wymiar dwóch swobód dla przepływu mieszanego potwierdza dokładny
układ wszystkich współczynników kwadratowych dla $n=3$: rzędy
$71/81$ bez wzajemności oraz $79/81$ po jej nałożeniu.
Nie jest to zastępstwo dowodu ogólnego. Gęsty racjonalny przykład
daje dokładną sprzeczność $\operatorname{rank}M=9$,
$\operatorname{rank}[M|b]=10$ nawet dla pełnego hermitowskiego
bloku antysymetrycznego. Przeciwne chiralności rzędu sześć,
zmiany niewidocznego bloku i pola lokalne na grafie są niezależnymi
kontrolami formuły \eqref{eq:entry}. Małe reszty numeryczne nie są
podstawą granicy \eqref{eq:number}.

\paragraph{Co pozostaje warunkowe.}
Klasa dotyczy stałych hermitowskich interakcji dwuciałowych
symetrycznych względem wymiany, z określonym czystym przepływem.
Pamięć kontrolera, inne prawa składania, otwarte układy i inne
mapy źródłowe nie zostały wykluczone. Protokół świeżych kopii
z poprzedniego raportu nie jest równowagą, więc ten wynik go nie obala.
Podstawowe techniki map komutujących, polaryzacji i nierówności
entropowych są znane; nie zgłasza się dla nich pierwszeństwa.
Nową treścią dla badanego FIN jest połączenie klasyfikacji z
niekasowalnym świadkiem strict i jego jawną granicą korelacyjną.

\paragraph{Najbardziej rozstrzygający dalszy obiekt.}
Potrzebny jest konkretny, niezależnie uzasadniony stan lub mechanizm
tworzenia korelacji $\chi$, który ma zerowe marginały, zachowuje
dodatniość $C\otimes E+\chi$ i spełnia
$[H,\chi]=-[H,C\otimes E]$. Musi w szczególności opłacić świadek
$\Tr(\chi B)=-\kappa\Delta$ oraz granicę \eqref{eq:floor},
a następnie przewidzieć mierzalną propagację i poprawki.
Samo dopasowanie jądra albo zmiana niewidocznego bloku nie realizuje
tego zadania. To jest konkretny nowy warunek źródłowy, nie kolejny
ogólny audyt zamkniętej ścieżki.

Nie wyprowadzono jednostek, fizycznego przygotowania ani zapisu
laboratoryjnego. Nie domknięto QW-2191, mostu legacy--strict,
transferu ról, SM/GR, $L_{\rm total}$ ani ToE. Nadsoliton pozostaje
pierwotną informacją w stanie solitonicznym; dodatkowa warstwa
informacyjna pod nim nie została założona. Strukturalne testy legacy
nie przenoszą na niego liczby \eqref{eq:number} ani fizycznych ról strict.

\section*{Źródła i dane}
\begin{enumerate}
\item Dokumenty źródłowe repozytorium wymienione w przypisach:
proponowany układ stanu--jądra, granica ST8648 i poprzedni raport
o kwantowym źródle. Są przesłankami i kontekstem, nie nowymi odkryciami.
\item M.~Brešar, Š.~Špenko, \emph{Functional identities in one variable},
§3, \href{https://users.fmf.uni-lj.si/bresar/documents/FIonevar.pdf}{wersja autorska}.
\item E.~A.~Carlen, E.~H.~Lieb, \emph{Remainder Terms for Some Quantum
Entropy Inequalities} (2014),
\href{https://arxiv.org/abs/1402.3840}{arXiv:1402.3840}.
\item Dokładne przedziały strict:
\texttt{fin\_replication\_consistency/certify.py} i
\texttt{fin\_projected\_learning/research.py}.
Kanoniczna referencja legacy:
\texttt{fin\_projected\_learning/geometry.py}.
\item Dane obecnego opracowania:
\texttt{fin\_hartree\_equivalence/research.py},
\texttt{results.json}, \texttt{test\_research.py} oraz
\texttt{verification.json}. Dowody ogólne, dokładne certyfikaty
skończone i kontrole zmiennoprzecinkowe mają odrębne role.
\end{enumerate}
\end{document}
